\documentclass[10pt]{article}
\usepackage[margin=1in]{geometry}
\usepackage{amsmath, amssymb}
\usepackage{enumitem}

\title{Lesson Plan: The Remainder Theorem}
\author{}
\date{}

\begin{document}

\maketitle

\section*{Objective}
Students will understand and apply the Remainder Theorem to evaluate polynomials and determine remainders without performing long division.

\section*{Standards}
\begin{itemize}
    \item Understand the relationship between polynomial division and function evaluation.
    \item Apply algebraic reasoning to simplify computations.
\end{itemize}

\section*{Prerequisites}
\begin{itemize}
    \item Polynomial evaluation
    \item Synthetic division (optional but helpful)
    \item Basic understanding of division of polynomials
\end{itemize}

\section*{Materials}
\begin{itemize}
    \item Whiteboard / projector
    \item Student worksheet
    \item Graphing calculator (optional)
\end{itemize}

\section*{Lesson Duration}
45--60 minutes

\section*{Lesson Structure}

\subsection*{1. Warm-Up (5--10 minutes)}
Evaluate the following:
\[
f(x) = x^2 + 3x + 2
\]
\begin{itemize}
    \item Find $f(1)$
    \item Find $f(-2)$
\end{itemize}

Prompt discussion: What does it mean to ``plug in'' a value?

\subsection*{2. Motivation (10 minutes)}
Ask students:
\begin{quote}
What if I divide $f(x) = x^2 + 3x + 2$ by $(x - 1)$? What remainder do I get?
\end{quote}

Walk through polynomial division (or synthetic division) to show:
\[
\text{Remainder} = 6
\]

Then connect:
\[
f(1) = 6
\]

\textbf{Key Question:} Is this a coincidence?

\subsection*{3. Direct Instruction (10--15 minutes)}
State the theorem clearly:

\begin{quote}
\textbf{Remainder Theorem:}  
If a polynomial $f(x)$ is divided by $(x - c)$, then the remainder is $f(c)$.
\end{quote}

\textbf{Example 1:}
\[
f(x) = 2x^3 - x + 5
\]
Find the remainder when divided by $(x - 2)$.

\[
f(2) = 2(2)^3 - 2 + 5 = 16 - 2 + 5 = 19
\]

\textbf{Conclusion:} The remainder is 19.

\subsection*{4. Guided Practice (10--15 minutes)}

Students try:

\begin{enumerate}[label=\arabic*.]
    \item $f(x) = x^3 - 4x + 1$, divided by $(x - 3)$
    \item $f(x) = 2x^2 + 5x - 7$, divided by $(x + 1)$
    \item $f(x) = x^4 - x^2 + 3$, divided by $(x - 0)$
\end{enumerate}

Encourage students to evaluate $f(c)$ instead of dividing.

\subsection*{5. Independent Practice (10 minutes)}

\begin{enumerate}[label=\arabic*.]
    \item Find the remainder when $f(x) = 3x^3 - 2x^2 + x - 5$ is divided by $(x - 1)$.
    \item Determine the remainder when $f(x) = x^4 + 2x^3 - x + 6$ is divided by $(x + 2)$.
    \item Challenge: Find $k$ such that the remainder when $f(x) = x^2 + kx + 4$ is divided by $(x - 2)$ is 10.
\end{enumerate}

\subsection*{6. Closure (5 minutes)}

Ask:
\begin{itemize}
    \item What does $f(c)$ represent in terms of division?
    \item Why is this theorem useful?
\end{itemize}

Summarize:
\begin{quote}
The Remainder Theorem allows us to avoid long division by simply evaluating the polynomial.
\end{quote}

\section*{Assessment}
\begin{itemize}
    \item Exit Ticket:
    \[
    f(x) = 2x^3 + x - 4
    \]
    Find the remainder when divided by $(x - 3)$.
\end{itemize}

\section*{Differentiation}

\textbf{Support:}
\begin{itemize}
    \item Provide step-by-step substitution examples
    \item Allow use of calculators
\end{itemize}

\textbf{Extension:}
\begin{itemize}
    \item Introduce Factor Theorem
    \item Connect to graphing: zeros and x-intercepts
\end{itemize}

\end{document}